\begin{frame}{self-play가 효과적인 이유}
  \begin{itemize}
    \item \textbf{(1) 난이도가 학습자와 \textbf{동행}} ---
      ``너무 쉬워서 배울 것이 없어지거나'' ``너무 어려워서
      신호가 없어지는'' 극단이 \textbf{없음}
    \item \textbf{(2) ``나보다 약간 강한 상대''를 이기려는
      학습}은 \textbf{정밀하게 차이만 교정}하는 신호를 만들어
      학습 효율이 \textbf{좋음} (15.2절 MCTS가 신경망보다
      정확한 수를 찾아줬던 것과 같은 이치)
  \end{itemize}
  \vskip0.6em
  \begin{itemize}
    \item Week 5의 몬테카를로 학습처럼 ``무작위 행동으로
      시작해 경험을 쌓아야 하는'' 학습에서는 \textbf{초기
      신호의 품질}이 병목 --- self-play는 ``\textbf{초기부터
      좋은 상대를 대결시킨}'' 버전
  \end{itemize}
  \vskip0.5em
  \begin{block}{다만}
    이 ``좋은 상대''를 사람이 만들어준 것(인간 게임
    데이터, AlphaGo(2016)의 경우)이 아니라 \textbf{학습
    과정 자체가 만들어내는} 것이 \textbf{AlphaZero(2017)의
    특징}.
  \end{block}
\end{frame}
